\chapter{Validation and Testing}

\section{Testing Methodology}
Testing was performed using functional testing, integration testing, UI testing, API testing, and log-based result validation. The focus was to verify that the system accepts images, returns structured responses, handles invalid model output safely, maps risk to correct actions, and persists records.

The testing flow was:
\begin{enumerate}
\item Upload normal landmark image.
\item Verify low-risk response and learn-more action.
\item Upload hazardous image.
\item Verify high-risk response and report action.
\item Test ambiguous or invalid response behavior.
\item Verify database records for messages, outputs, actions, and incidents.
\item Verify report confirmation and optional notification status.
\end{enumerate}

\section{Test Cases}
\begin{table}[H]
\centering
\caption{Test Case Summary}
\renewcommand{\arraystretch}{1.3}
\small
\begin{tabularx}{\textwidth}{|c|X|X|c|X|}
\hline
\textbf{Sr. No} & \textbf{Test Case Title / Description} & \textbf{Testing Methodology Applied} & \textbf{Result} & \textbf{Defects Identified} \\
\hline
1 & Upload supported image file and submit for analysis & Functional testing & PASS & None \\
\hline
2 & Analyze regular Shaniwar Wada image & System testing & PASS & None \\
\hline
3 & Analyze fire-affected Shaniwar Wada image & System testing & PASS & None \\
\hline
4 & Verify high-risk output shows Report action & Black-box testing & PASS & None \\
\hline
5 & Verify low-risk output shows Learn More and Open Map actions & Black-box testing & PASS & None \\
\hline
6 & Verify uncertain or low-confidence response asks for more context & Boundary testing & PASS & None \\
\hline
7 & Verify report action requires confirmation before incident creation & Functional testing & PASS & None \\
\hline
8 & Verify model output, action, and incident records are saved in SQLite & Integration testing & PASS & None \\
\hline
9 & Verify missing/invalid model response returns safe fallback & Negative testing & PASS & None \\
\hline
10 & Verify learn-more action fetches external knowledge summary & Integration testing & PASS & Depends on Wikipedia availability \\
\hline
\end{tabularx}
\end{table}

\section{Validation of Results}
The stored logs show 22 analyzed runs across 14 sessions. The risk distribution includes 8 low-risk outputs, 8 high-risk outputs, and 6 uncertain outputs. Confidence values range from 0.00 to 0.95 with an average of 0.623.

\begin{table}[H]
\centering
\caption{Confidence Band Distribution}
\begin{tabular}{|l|c|c|}
\hline
\textbf{Band} & \textbf{Count} & \textbf{Share} \\ \hline
Low ($<0.30$) & 6 & 27.3\% \\ \hline
Mid ($0.30$--$0.79$) & 5 & 22.7\% \\ \hline
High ($\geq0.80$) & 11 & 50.0\% \\ \hline
\end{tabular}
\end{table}

The case-study images validate the main intended behavior: the normal image leads to an informational action, while the hazardous image leads to a report-oriented action. The system therefore satisfies the MVP acceptance criteria for normal image interpretation, risky image escalation, ambiguity handling, and incident log review.
